% !TEX root = ../../main.tex

\subsection{Design Insights and Chapter Summary}

This chapter evaluated the baseline OFDM-LMMSE, baseline OTFS, and OTFS-IM systems under the same delay-Doppler channel and modulation assumptions. The results show that OTFS-IM is not simply an additional mapping layer placed on top of OTFS. Its performance depends jointly on the IM block size, the number of active positions, the spectral efficiency, the selected activation-pattern table, and the reliability of the detector output. Therefore, the most meaningful conclusion is obtained by reading the BER curves together with the error-decomposition and BER-SE trade-off results.

The first important observation is that OTFS-IM can improve BER over baseline OTFS in the simulated high-Doppler delay-Doppler channel. This is most clearly shown in the medium-to-high-SNR region, where the detector output becomes reliable enough for the activation pattern and symbol decisions to separate more cleanly. In the equal-spectral-efficiency cases, especially \(n=6,k=5\), OTFS-IM preserves the same nominal rate as baseline OTFS while reducing BER. This is the strongest result in the chapter because it shows that the gain does not only come from lowering the number of transmitted information bits.

At low SNR, the advantage of OTFS-IM is less consistent. In this region, the received samples are strongly affected by noise, and the MP detector has difficulty distinguishing both the active positions and the transmitted QAM symbols. As a result, the total BER is influenced more by detector uncertainty than by the exact geometry of the selected pattern table. This explains why OHD, B-OHD, and random selection can be close to each other at some SNR points. It also explains why pattern selection alone cannot be expected to solve all low-SNR errors.

The error-decomposition results provide a more detailed view of the OTFS-IM error mechanism. The total BER is produced by two different sources: index-bit errors and symbol-bit errors. Index errors occur when the receiver selects the wrong active-position pattern, while symbol errors occur when the active positions are detected but the QAM symbols are decoded incorrectly. In the simulated results, symbol errors remain a major contribution to the total BER, especially at moderate SNR. This means that improving pattern geometry is useful, but it should be combined with reliable detection and symbol estimation. A pattern table with good Hamming-distance properties cannot fully compensate for poor detector reliability.

The comparison between OHD and B-OHD gives a more careful interpretation of the proposed pattern-selection improvement. OHD selects activation patterns by prioritizing Hamming-distance separation. B-OHD keeps this idea but additionally considers carrier-usage balance. This makes B-OHD a refinement of OHD rather than a different receiver or a different modulation scheme. The gain of B-OHD appears when the OHD table is unbalanced and when the selected table can be modified without reducing the minimum Hamming distance. The \(n=5,k=3\) configuration illustrates this behavior because B-OHD improves the usage distribution and provides better performance than OHD at the selected operating points.

However, B-OHD does not always produce a different pattern table. In configurations such as \(n=6,k=5\), the selected OHD and B-OHD tables can become effectively identical, leading to overlapping BER curves. This is not a simulation error. It means that the OHD solution is already balanced enough under that configuration, or that the available candidate patterns do not allow B-OHD to produce a meaningfully different table. This result is useful because it prevents an overclaimed conclusion: B-OHD is beneficial when the pattern-selection problem has room for balancing, but it is not universally superior in every \((n,k)\) setting.

The random-pattern baseline also plays an important role in the evaluation. It shows whether a deterministic pattern-selection rule is truly useful or whether similar performance could be obtained from an arbitrary pattern table. In some cases, the random average is competitive with OHD or B-OHD, particularly when the number of candidate patterns is small or when several pattern tables have similar distance properties. Therefore, the random baseline should not be removed from the analysis. Its presence makes the evaluation more honest and strengthens the conclusion that pattern selection is configuration-dependent.

The BER-SE trade-off analysis further clarifies how the configurations should be interpreted. Equal-SE configurations such as \(n=4,k=3\), \(n=5,k=4\), and \(n=6,k=5\) provide the fairest direct comparison with baseline OTFS because they preserve the same spectral efficiency of \(2.00~\text{bits/symbol}\). Lower-SE configurations, such as \(n=4,k=2\), \(n=5,k=3\), \(n=6,k=3\), and \(n=6,k=4\), may achieve lower BER, but part of that improvement comes from transmitting fewer information bits per IM block. These lower-SE configurations are still useful because practical communication systems may choose reliability over rate in difficult channel conditions, but their results should be described as a rate-reliability trade-off rather than a strict performance gain at equal rate.

From a design perspective, the simulation results suggest that the choice of \((n,k)\) should not be made only from the spectral-efficiency formula. A larger \(k\) can help preserve rate, while a smaller \(k\) may improve reliability by reducing the number of active symbols and changing the pattern-detection problem. At the same time, increasing \(n\) enlarges the pattern space and can make exact pattern selection more expensive. Therefore, the final OTFS-IM design must balance spectral efficiency, BER performance, pattern-selection complexity, and detector reliability.

The results also reveal several limitations of the present simulation framework. First, the channel is modeled in the delay-Doppler domain with discrete delay and Doppler taps, but the simulation does not sweep a physical velocity parameter. Therefore, the results support the evaluation of a high-Doppler or doubly selective channel model, while a separate velocity-based study would be needed to quantify performance at specific speeds such as vehicular, railway, or UAV scenarios. Second, the Monte Carlo simulation uses a finite number of frames, so very low BER points may still contain statistical variation. Third, the receiver uses the same general MP detection principle for the compared OTFS-IM designs; more advanced channel-aware, detector-aware, or soft-output detection methods may further change the relative performance.

In summary, this chapter confirms that OTFS-IM is a meaningful extension of baseline OTFS for delay-Doppler-domain transmission. The simulations show that OTFS-IM can improve BER, including in equal-SE cases, and that pattern-table design affects the reliability of index detection. The proposed B-OHD method improves OHD when carrier-usage balancing changes the selected pattern table, but its gain is configuration-dependent. These findings provide the basis for the final conclusion of the thesis: OTFS-IM offers a flexible BER-SE design space for high-mobility-oriented wireless channels, and structured pattern selection is a useful but not standalone component of that design.
